%% kapa: exam=1 round=27 year=2016 period=1 form=A subject=civil_law subject_ko=민법 seq=31 number=31
%% source: Q-Net 감정평가사 기출문제 (cst00309) artlSeq=5169402 file=extracted/1008/5169402/제27회 감정평가사 1차 문제/제27회 감정평가사 1교시문제(A형).pdf
%% conversion: kordoc 4.24.51 -> tools/gen1.py (auto); review=proofread-sonnet; flags=none
\begin{question}{31}
점유자와 회복자의 법률관계에 관한 설명으로 옳지 않은 것은? (다툼이 있으면 판례에 따름)
\begin{choices}
  \item 타인의 건물을 선의로 점유한 자는 비록 법률상 원인 없이 사용하였더라도 이로 인한 이득을 반환할 의무가 없다.
  \item 악의의 점유자가 과실을 소비한 경우에는 그 과실의 대가를 보상하여야 한다.
  \item 점유물이 점유자의 책임 있는 사유로 인하여 멸실 또는 훼손된 경우, 선의의 자주점유자는 그 이익이 현존하는 한도에서 배상하여야 한다.
  \item 선의의 점유자가 본권에 관한 소에서 패소한 경우, 제소 후 판결확정 전에 취득한 과실은 반환할 의무가 없다.
  \item 점유자가 과실을 취득한 경우에는 통상의 필요비의 상환을 청구하지 못한다.
\end{choices}
\end{question}
\answer{4}
